\documentclass[12pt,a4paper]{article}

\usepackage[utf8]{vietnam}
\usepackage{amsmath}
\usepackage{amssymb}
\usepackage{amsfonts}
\usepackage{geometry}
\usepackage{longtable}
\usepackage{array}
\usepackage{hyperref}
\geometry{a4paper, margin=1in}

\title{Phân tích thuật toán xử lý QR Code bằng Camera và LiDAR}
\author{Robot Dogzilla}
\date{\today}

\begin{document}

\maketitle

\section{Mục tiêu}

Hệ thống hiện tại xử lý QR theo hướng:

\begin{itemize}
    \item Camera dùng để phát hiện QR Code, lấy 4 góc ảnh, ước lượng pose sơ bộ và suy ra góc lệch.
    \item LiDAR dùng để hiệu chỉnh khoảng cách thực theo hướng QR khi tìm được điểm scan phù hợp trên ray.
    \item Overlay hiển thị khoảng cách theo LiDAR nếu có, nếu không thì fallback về khoảng cách camera.
\end{itemize}

Kết luận quan trọng:

\begin{itemize}
    \item \texttt{distance\_m} ban đầu từ camera chỉ là giá trị sơ bộ.
    \item Khi có hit LiDAR hợp lệ, \texttt{distance\_m}, \texttt{lateral\_x\_m}, \texttt{forward\_z\_m}, \texttt{target\_*\_m} được ghi đè theo LiDAR.
    \item \texttt{overlay.py} chỉ hiển thị giá trị camera khi \texttt{lidar\_distance\_m} không hợp lệ.
\end{itemize}

% ============================================================
\section{Luồng xử lý tổng quát}

Pipeline thực tế trong code:

\begin{enumerate}
    \item Đọc frame từ camera robot.
    \item Phát hiện QR bằng \texttt{pyzbar} hoặc fallback OpenCV.
    \item Sắp xếp 4 góc QR theo thứ tự top-left, top-right, bottom-right, bottom-left.
    \item Chạy \texttt{solvePnP} để lấy pose sơ bộ.
    \item Tính góc QR và khoảng cách camera.
    \item Truy vấn LiDAR theo bearing của QR trên map.
    \item Nếu tìm được hit LiDAR, thay thế toàn bộ phần range liên quan bằng giá trị LiDAR.
    \item Sinh payload và hiển thị overlay.
\end{enumerate}

% ============================================================
\section{Camera và ma trận nội tại}

Trong \texttt{qr\_detect.py}, ma trận camera được khai báo:

\begin{equation}
K =
\begin{bmatrix}
920 & 0 & 640 \\
0 & 920 & 360 \\
0 & 0 & 1
\end{bmatrix}
\end{equation}

Ý nghĩa:

\begin{longtable}{|>{\raggedright\arraybackslash}p{4cm}|>{\raggedright\arraybackslash}p{9cm}|}
\hline
Biến & Ý nghĩa \\ \hline
$f_x = 920$ & tiêu cự theo trục $X$ \\ \hline
$f_y = 920$ & tiêu cự theo trục $Y$ \\ \hline
$c_x = 640$ & tâm ảnh theo trục $X$ \\ \hline
$c_y = 360$ & tâm ảnh theo trục $Y$ \\ \hline
\end{longtable}

Hệ tọa độ camera trong phần tính pose:

\begin{itemize}
    \item Trục $X$: ngang trái/phải.
    \item Trục $Y$: dọc lên/xuống.
    \item Trục $Z$: hướng nhìn tiến về phía trước camera.
\end{itemize}

% ============================================================
\section{Tiền xử lý và phát hiện QR}

Frame được nén xuống độ rộng tối đa \texttt{detect\_width = 640} trước khi decode QR:

\begin{equation}
scale = \frac{detect\_width}{w}
\end{equation}

Nếu frame có chiều rộng $w \le detect\_width$ thì:

\begin{equation}
scale = 1
\end{equation}

Sau đó ảnh được chuyển grayscale và làm mượt nhẹ bằng Gaussian blur.

Điểm quan trọng là các điểm ảnh QR sau khi decode được đưa ngược về kích thước frame gốc bằng:

\begin{equation}
P_{img}^{(orig)} = \frac{P_{img}^{(scaled)}}{scale}
\end{equation}

% ============================================================
\section{Sắp xếp 4 góc QR}

Các góc QR được sắp theo thứ tự:

\begin{enumerate}
    \item top-left
    \item top-right
    \item bottom-right
    \item bottom-left
\end{enumerate}

Giả sử 4 điểm ảnh là:

\begin{equation}
P_{img} =
\begin{bmatrix}
u_1 & v_1 \\
u_2 & v_2 \\
u_3 & v_3 \\
u_4 & v_4
\end{bmatrix}
\end{equation}

Khi cần fallback từ polygon sang bounding box, code dùng hình chữ nhật bao quanh.

% ============================================================
\section{Object points của QR}

Kích thước thật của QR:

\begin{equation}
S = 0.12 \text{ m}
\end{equation}

Đặt:

\begin{equation}
s = \frac{S}{2}
\end{equation}

Bốn điểm 3D của QR trên mặt phẳng $Z = 0$:

\begin{equation}
P_{obj} =
\begin{bmatrix}
-s & -s & 0 \\
 s & -s & 0 \\
 s &  s & 0 \\
-s &  s & 0
\end{bmatrix}
\end{equation}

Ý nghĩa:

\begin{itemize}
    \item Tâm QR là gốc tọa độ local của object.
    \item Mặt phẳng QR được giả định song song mặt phẳng $XY$ của object frame.
\end{itemize}

% ============================================================
\section{Ước lượng pose bằng \texttt{solvePnP}}

Hàm \texttt{cv2.solvePnP} trả về:

\begin{equation}
(rvec, tvec)
\end{equation}

Trong đó:

\begin{equation}
tvec =
\begin{bmatrix}
t_x \\
t_y \\
t_z
\end{bmatrix}
\end{equation}

Code dùng:

\begin{itemize}
    \item \texttt{SOLVEPNP\_ITERATIVE} trước.
    \item Nếu thất bại thì thử \texttt{SOLVEPNP\_IPPE\_SQUARE}.
\end{itemize}

Giá trị trục tiến trước được chuẩn hóa bằng:

\begin{equation}
tz = |t_z|
\end{equation}

Lưu ý:

\begin{itemize}
    \item Code không dùng trực tiếp $t_x$ từ \texttt{solvePnP}.
    \item Tọa độ ngang được tính lại từ tâm ảnh bằng công thức pinhole.
\end{itemize}

% ============================================================
\section{Tính lệch ngang từ tâm ảnh}

Tâm QR trên ảnh được tính bằng:

\begin{equation}
u_c = \frac{u_1+u_2+u_3+u_4}{4}
\end{equation}

\begin{equation}
v_c = \frac{v_1+v_2+v_3+v_4}{4}
\end{equation}

Trong code, lệch ngang được tính theo mô hình pinhole:

\begin{equation}
t_x = \frac{u_c - c_x'}{f_x'} \cdot t_z
\end{equation}

trong đó $f_x'$ và $c_x'$ là phần tử sau khi scale ma trận camera theo frame thực tế.

Hệ quả:

\begin{itemize}
    \item $t_x > 0$: QR lệch sang phải trong hệ camera theo quy ước hiện tại.
    \item $t_x < 0$: QR lệch sang trái.
\end{itemize}

% ============================================================
\section{Góc QR}

Góc lệch của QR so với trục nhìn được tính bằng:

\begin{equation}
\theta = \operatorname{atan2}(t_x, t_z)
\end{equation}

Đổi sang độ:

\begin{equation}
\theta_{deg} = \theta \cdot \frac{180}{\pi}
\end{equation}

Phân loại hướng:

\begin{equation}
direction =
\begin{cases}
\text{right}, & \theta_{deg} > 5^\circ \\
\text{left}, & \theta_{deg} < -5^\circ \\
\text{center}, & \text{ngược lại}
\end{cases}
\end{equation}

% ============================================================
\section{Khoảng cách camera sơ bộ}

Khoảng cách từ camera đến QR được tính như độ dài vector trong hệ camera:

\begin{equation}
d_{camera} = \sqrt{t_x^2 + t_z^2}
\end{equation}

Giá trị này được gán vào:

\begin{itemize}
    \item \texttt{distance\_m}
    \item \texttt{camera\_distance\_m}
\end{itemize}

và \texttt{lidar\_distance\_m} ban đầu để \texttt{None}.

% ============================================================
\section{Tọa độ target theo camera}

Nếu chỉ dùng camera, target được suy ra trực tiếp từ pose camera:

\begin{equation}
target_x = t_x
\end{equation}

\begin{equation}
target_z = t_z
\end{equation}

\begin{equation}
target_{distance} = \sqrt{t_x^2 + t_z^2}
\end{equation}

Nói cách khác, target ban đầu chính là vector tương đối của QR trong hệ robot.

% ============================================================
\section{Truy vấn LiDAR theo bearing QR}

Phần LiDAR được xử lý trong \texttt{find\_first\_obstacle\_on\_ray()}.

Góc tia quét:

\begin{equation}
\phi = \operatorname{normalize}(yaw_{robot} + \theta + offset)
\end{equation}

trong đó:

\begin{itemize}
    \item $yaw_{robot}$ là yaw robot trong map frame.
    \item $\theta$ là góc QR trong camera frame.
    \item $offset$ là \texttt{CAMERA\_LIDAR\_YAW\_OFFSET\_RAD}.
\end{itemize}

Với mỗi điểm scan trên map:

\begin{equation}
dx = p_x - r_x,\quad dy = p_y - r_y
\end{equation}

Khoảng cách chiếu theo tia:

\begin{equation}
forward_m = dx \cos(\phi) + dy \sin(\phi)
\end{equation}

Khoảng cách vuông góc đến tia:

\begin{equation}
perp_m = | -dx \sin(\phi) + dy \cos(\phi) |
\end{equation}

Điểm scan hợp lệ phải thỏa:

\begin{equation}
forward_m \ge min\_distance\_m
\end{equation}

\begin{equation}
perp_m \le corridor\_width\_m
\end{equation}

Trong các điểm hợp lệ, điểm có $forward_m$ nhỏ nhất sẽ được chọn làm hit đầu tiên trên ray.

Kết quả trả về:

\begin{equation}
ray\_distance\_m = forward_m
\end{equation}

\begin{equation}
ray\_angle\_rad = \phi
\end{equation}

% ============================================================
\section{Hiệu chỉnh bằng LiDAR}

Khi có hit LiDAR hợp lệ, code ghi đè các giá trị liên quan đến range:

\begin{equation}
distance_m = ray\_distance\_m
\end{equation}

\begin{equation}
lateral\_x\_m = ray\_distance\_m \sin(ray\_angle\_rad)
\end{equation}

\begin{equation}
forward\_z\_m = ray\_distance\_m \cos(ray\_angle\_rad)
\end{equation}

\begin{equation}
target_x\_m = lateral\_x\_m
\end{equation}

\begin{equation}
target_z\_m = forward\_z\_m
\end{equation}

\begin{equation}
target\_distance\_m = ray\_distance\_m
\end{equation}

\begin{equation}
distance\_source = \text{"lidar"}
\end{equation}

Ngoài ra:

\begin{itemize}
    \item \texttt{map\_x\_m}, \texttt{map\_y\_m} lưu tọa độ điểm hit trên map.
    \item \texttt{ray\_distance\_m} là khoảng cách dọc theo tia ray.
    \item \texttt{distance\_m} sau enrichment dùng đúng giá trị LiDAR.
\end{itemize}

% ============================================================
\section{Ý nghĩa các trường dữ liệu}

\begin{longtable}{|>{\raggedright\arraybackslash}p{4cm}|>{\raggedright\arraybackslash}p{9cm}|}
\hline
Trường & Ý nghĩa \\ \hline
\texttt{angle\_deg} & góc lệch QR, đơn vị độ \\ \hline
\texttt{angle\_rad} & góc lệch QR, đơn vị radian \\ \hline
\texttt{distance\_m} & khoảng cách đang dùng cho xử lý, camera hoặc LiDAR tùy nguồn \\ \hline
\texttt{camera\_distance\_m} & khoảng cách sơ bộ từ camera \\ \hline
\texttt{lidar\_distance\_m} & khoảng cách từ LiDAR nếu có hit hợp lệ \\ \hline
\texttt{lateral\_x\_m} & thành phần ngang trong hệ robot \\ \hline
\texttt{forward\_z\_m} & thành phần tiến trước trong hệ robot \\ \hline
\texttt{target\_x\_m} & hoành độ target trong hệ robot \\ \hline
\texttt{target\_z\_m} & tung độ target trong hệ robot \\ \hline
\texttt{target\_distance\_m} & khoảng cách đến target đang sử dụng \\ \hline
\texttt{ray\_distance\_m} & khoảng cách dọc theo ray LiDAR \\ \hline
\texttt{ray\_angle\_rad} & góc tia LiDAR trong map frame \\ \hline
\texttt{distance\_source} & nguồn khoảng cách: \texttt{lidar} hoặc \texttt{camera} \\ \hline
\end{longtable}

% ============================================================
\section{Hiển thị overlay}

Trong \texttt{overlay.py}, khoảng cách hiển thị được chọn theo quy tắc:

\begin{equation}
display\_distance =
\begin{cases}
lidar\_distance\_m, & \text{nếu hợp lệ và finite} \\
distance\_m, & \text{ngược lại}
\end{cases}
\end{equation}

Chuỗi hiển thị chính:

\begin{verbatim}
QR: ...
angle: ... deg
dist : ... m
tx/t: (..., ...)
target: (..., ...)
map: (..., ...)
ray: ... m
\end{verbatim}

Trong đó:

\begin{itemize}
    \item \texttt{dist} là khoảng cách nguồn đang dùng để hiển thị.
    \item \texttt{tx/t} hiển thị \texttt{lateral\_x\_m} và \texttt{display\_distance}.
    \item \texttt{target} hiển thị \texttt{target\_x\_m} và \texttt{target\_z\_m}.
    \item \texttt{ray} chỉ xuất hiện khi có \texttt{ray\_distance\_m}.
\end{itemize}

% ============================================================
\section{Liên hệ với code}

\begin{itemize}
    \item \texttt{backend/control/services/qr\_detector.py}
    \begin{itemize}
        \item decode QR,
        \item sắp xếp 4 góc,
        \item chạy \texttt{solvePnP},
        \item tính \texttt{angle\_rad}, \texttt{angle\_deg}, \texttt{distance\_m} sơ bộ.
    \end{itemize}
    \item \texttt{backend/control/services/qr\_detect.py}
    \begin{itemize}
        \item gọi detector,
        \item truy vấn SLAM/LiDAR,
        \item enrichment theo ray,
        \item tạo payload xuất ra UI.
    \end{itemize}
    \item \texttt{backend/control/services/overlay.py}
    \begin{itemize}
        \item chọn nguồn khoảng cách để hiển thị,
        \item render text overlay theo dữ liệu đã enrich.
    \end{itemize}
\end{itemize}

% ============================================================
\section{Kết luận}

Mô hình xử lý QR hiện tại có thể tóm tắt bằng chuỗi:

\begin{equation}
Image \rightarrow QR\ Detection \rightarrow Pose \rightarrow Bearing \rightarrow LiDAR\ Raycast \rightarrow Range\ Override \rightarrow Target \rightarrow Overlay
\end{equation}

Tóm lược ngắn:

\begin{itemize}
    \item Camera chịu trách nhiệm nhận diện QR và suy ra góc.
    \item Camera chỉ cung cấp khoảng cách sơ bộ trước khi LiDAR can thiệp.
    \item LiDAR quyết định khoảng cách thực khi tìm được hit hợp lệ.
    \item Overlay phản ánh đúng nguồn đang được dùng.
\end{itemize}

\end{document}
